% ----------------------------------------------------------------------
% Section 2: Conventions -- spinors, superspace and superfields
% Writer: Writer_Notation. Sources: extract-01-spinor-superspace.md (primary),
% extract-02-action-bv.md [only for (3A.31)-(3A.32), same repo source step-03a],
% extract-03-extended-grammar.md (notation dictionary cross-check).
% Pure body LaTeX per NOTATION.md: one \section, no \newcommand.
% ----------------------------------------------------------------------
\section{Conventions: spinors, superspace and superfields}\label{sec:notation}

We use two-component Weyl spinors throughout, in both Lorentzian and Euclidean
signature; textbook references for the superspace conventions are
\cite{GGRS,WeinbergIII}. Lorentzian vector indices are $\mu,\nu,\rho=0,1,2,3$,
Euclidean vector indices are $m,n,r=1,2,3,4$, undotted spinor indices are
$\alpha,\beta,\gamma=1,2$, and dotted spinor indices are
$\dot\alpha,\dot\beta,\dot\gamma=\dot1,\dot2$. The index alphabet already
distinguishes the two signatures; where confusion is possible we add the label
$L$ (Lorentzian) or $E$ (Euclidean) explicitly. All commutators are graded,
$[A,B\}:=AB-(-1)^{|A||B|}BA$, with $|X|\in\{0,1\}$ the Grassmann parity.

\subsection{Lorentzian spinors}\label{sec:notation:lorentzian}

We use the mostly-plus metric and the orientation
\begin{equation}
\eta_{\mu\nu}=\operatorname{diag}(-1,+1,+1,+1),\qquad
v_\mu=\eta_{\mu\nu}v^\nu,\qquad
\epsilon^{0123}=+1 .
\label{eq:2:metric}
\end{equation}
The two-component epsilon tensors and index raising and lowering are
\begin{equation}
\epsilon^{12}=\epsilon^{\dot1\dot2}=+1,\qquad
\epsilon_{12}=\epsilon_{\dot1\dot2}=-1,\qquad
\psi^\alpha=\epsilon^{\alpha\beta}\psi_\beta,\qquad
\bar\psi^{\dot\alpha}=\epsilon^{\dot\alpha\dot\beta}\bar\psi_{\dot\beta},
\label{eq:2:epsilon}
\end{equation}
with $\epsilon^{\alpha\beta}\epsilon_{\beta\gamma}=\delta^\alpha{}_\gamma$,
$\epsilon_{\alpha\beta}\epsilon^{\beta\gamma}=\delta_\alpha{}^\gamma$, and the
same relations with dotted indices. Spinor contractions follow the
\emph{asymmetric} summation convention: undotted indices are contracted from
upper left to lower right, dotted indices from lower left to upper right,
\begin{equation}
\xi\chi:=\xi^\alpha\chi_\alpha,\qquad
\bar\xi\bar\chi:=\bar\xi_{\dot\alpha}\bar\chi^{\dot\alpha},\qquad
\xi^\alpha\chi_\alpha=-\,\xi_\alpha\chi^\alpha,\qquad
\bar\xi^{\dot\alpha}\bar\chi_{\dot\alpha}
=-\,\bar\xi_{\dot\alpha}\bar\chi^{\dot\alpha}.
\label{eq:2:contraction}
\end{equation}
For Grassmann-odd spinors this gives $\xi\chi=\chi\xi$ and
$\bar\xi\bar\chi=\bar\chi\bar\xi$. With $\sigma^1,\sigma^2,\sigma^3$ the
standard Pauli matrices, the sigma matrices are
\begin{equation}
(\sigma^\mu)_{\alpha\dot\alpha}
=(\mathbf 1_2,\sigma^1,\sigma^2,\sigma^3)_{\alpha\dot\alpha},\qquad
(\bar\sigma^\mu)^{\dot\alpha\alpha}
:=\epsilon^{\alpha\beta}\epsilon^{\dot\alpha\dot\beta}
(\sigma^\mu)_{\beta\dot\beta}
=(\mathbf 1_2,-\sigma^1,-\sigma^2,-\sigma^3)^{\dot\alpha\alpha},
\label{eq:2:sigmaL}
\end{equation}
and obey the Clifford relations
\begin{equation}
(\sigma^\mu)_{\alpha\dot\alpha}(\bar\sigma^\nu)^{\dot\alpha\beta}
+(\sigma^\nu)_{\alpha\dot\alpha}(\bar\sigma^\mu)^{\dot\alpha\beta}
=-2\eta^{\mu\nu}\delta_\alpha{}^\beta,\qquad
(\bar\sigma^\mu)^{\dot\alpha\alpha}(\sigma^\nu)_{\alpha\dot\beta}
+(\bar\sigma^\nu)^{\dot\alpha\alpha}(\sigma^\mu)_{\alpha\dot\beta}
=-2\eta^{\mu\nu}\delta^{\dot\alpha}{}_{\dot\beta}.
\label{eq:2:cliffordL}
\end{equation}
The (anti-)self-dual Lorentz generators are
\begin{equation}
(\sigma^{\mu\nu})_\alpha{}^\beta
:=\tfrac14\big[(\sigma^\mu)_{\alpha\dot\gamma}(\bar\sigma^\nu)^{\dot\gamma\beta}
-(\sigma^\nu)_{\alpha\dot\gamma}(\bar\sigma^\mu)^{\dot\gamma\beta}\big],\qquad
(\bar\sigma^{\mu\nu})^{\dot\alpha}{}_{\dot\beta}
:=\tfrac14\big[(\bar\sigma^\mu)^{\dot\alpha\gamma}(\sigma^\nu)_{\gamma\dot\beta}
-(\bar\sigma^\nu)^{\dot\alpha\gamma}(\sigma^\mu)_{\gamma\dot\beta}\big],
\label{eq:2:sigmamunu}
\end{equation}
with components $\sigma^{0i}=-\tfrac12\sigma^i$,
$\bar\sigma^{0i}=+\tfrac12\sigma^i$ and
$\sigma^{ij}=\bar\sigma^{ij}=-\tfrac{i}{2}\epsilon^{ijk}\sigma^k$.

\paragraph{Weyl versus Majorana spinors.} The formalism is strictly
two-component. In Lorentzian signature the barred spinor is the Hermitian
conjugate of the unbarred one,
\begin{equation}
\bar\chi^{\dot\alpha}:=(\chi^\alpha)^\dagger ;
\label{eq:2:majorana}
\end{equation}
this is the Majorana reality condition: a Majorana four-spinor is the pair
$(\chi_\alpha,\bar\chi^{\dot\alpha})$, while a Weyl spinor is half of this
data and is the fundamental object here. In particular the Lorentzian
supercharges obey $(Q_\alpha)^\dagger=\bar Q_{\dot\alpha}$.

\subsection{Euclidean spinors}\label{sec:notation:euclidean}

The Wick rotation keeps the spatial coordinates and the Grassmann coordinates
fixed and rotates the time coordinate; the generators acquire phases,
\begin{equation}
x_E^i=x_L^i,\qquad x_E^4=ix_L^0,\qquad
P^E_i=-iP^L_i,\qquad P^E_4=P_L^0=H,\qquad
Q^E_\alpha=-iQ^L_\alpha,\qquad
\bar Q^E_{\dot\alpha}=-i\bar Q^L_{\dot\alpha} .
\label{eq:2:wick}
\end{equation}
The Euclidean sigma matrices and Clifford relations are
\begin{equation}
(\sigma^m)_{\alpha\dot\alpha}
=(-i\sigma^1,-i\sigma^2,-i\sigma^3,\mathbf 1_2)_{\alpha\dot\alpha},\qquad
(\bar\sigma^m)^{\dot\alpha\alpha}
:=\epsilon^{\alpha\beta}\epsilon^{\dot\alpha\dot\beta}
(\sigma^m)_{\beta\dot\beta}
=(+i\sigma^1,+i\sigma^2,+i\sigma^3,\mathbf 1_2)^{\dot\alpha\alpha},
\label{eq:2:sigmaE}
\end{equation}
\begin{equation}
(\sigma^m)_{\alpha\dot\alpha}(\bar\sigma^n)^{\dot\alpha\beta}
+(\sigma^n)_{\alpha\dot\alpha}(\bar\sigma^m)^{\dot\alpha\beta}
=2\delta^{mn}\delta_\alpha{}^\beta ,
\label{eq:2:cliffordE}
\end{equation}
so that $(\bar\sigma^m)^{\dot\alpha\alpha}
=\big((\sigma^m)_{\alpha\dot\alpha}\big)^\dagger$ as matrices. The
(anti-)self-dual rotation generators $\sigma^{mn},\bar\sigma^{mn}$ are defined
as in \eqref{eq:2:sigmamunu} with $\mu\to m$, and have components
\begin{equation}
\sigma^{ij}=\bar\sigma^{ij}=\frac{i}{2}\epsilon^{ijk}\sigma^k,\qquad
\sigma^{4i}=+\frac{i}{2}\sigma^i,\qquad
\bar\sigma^{4i}=-\frac{i}{2}\sigma^i .
\label{eq:2:sigmamn}
\end{equation}

\paragraph{The $SO(4)\simeq SU(2)\times SU(2)$ structure.} In Euclidean
signature the undotted index $\alpha$ and the dotted index $\dot\alpha$ are
fundamental indices of two \emph{independent} $SU(2)$ factors:
$(\sigma^{mn})_\alpha{}^\beta$ acts only on undotted spinors and
$(\bar\sigma^{mn})^{\dot\alpha}{}_{\dot\beta}$ only on dotted ones. From
\eqref{eq:2:sigmamn}, $\sigma^{mn}$ and $\bar\sigma^{mn}$ are the
(anti-)self-dual generators of
$\mathfrak{so}(4)=\mathfrak{su}(2)_+\oplus\mathfrak{su}(2)_-$, interchanged by
$\sigma\leftrightarrow\bar\sigma$; the vector $(\sigma^m)_{\alpha\dot\alpha}$
is the $(\mathbf 2,\mathbf 2)$.

\paragraph{No Euclidean conjugation.} In Euclidean signature no dagger
relation is imposed between $Q_\alpha$ and $\bar Q_{\dot\alpha}$: the two Weyl
sectors are independent. Consequently the chiral and antichiral fields are
independent Euclidean fields, not conjugates. This holds for the matter
multiplets $(\Phi,\widetilde\Phi)$ and for the two gaugini
$(\lambda_\alpha,\widetilde\lambda_{\dot\alpha})$ alike, and it is used
throughout this paper: the superspace calculation of
Section~\ref{sec:superspace} and the component reduction of
Section~\ref{sec:component} treat untilded and tilded fields as independent
integration variables.
\begin{remark}\label{rem:2:noconjugation}
Standing convention: \emph{no Euclidean conjugation relates $\Phi$ to
$\widetilde\Phi$} (nor $\lambda$ to $\widetilde\lambda$). The Lorentzian
identifications $\widetilde\Phi=\bar\Phi$,
$\widetilde\lambda=\bar\lambda$ are Wick-transported relations only; no
intrinsic Euclidean reality condition is imposed anywhere in this paper.
\end{remark}

\paragraph{The plus/minus frame.} In the Euclidean computation of
Sections~\ref{sec:superspace}--\ref{sec:component} the undotted $SU(2)$ index
is resolved in the frame $\alpha=(+,-)$: an undotted spinor $X_\alpha$ has
components $X_\pm$, and correspondingly the gauge-covariant superspace
derivatives $\nabla_\alpha$ (undotted) and $\widetilde\nabla_{\dot\alpha}$
(dotted) have frame components
\begin{equation}
\nabla_\pm:=\nabla_{\alpha=\pm},\qquad
\widetilde\nabla_{\dot\pm}:=\widetilde\nabla_{\dot\alpha=\dot\pm} .
\label{eq:2:pmframe}
\end{equation}
Here $\nabla_\alpha$ carries the chiral gauge connection and
$\widetilde\nabla_{\dot\alpha}$ the antichiral one, adapted to the independent
Euclidean antichiral fields. In the one-loop calculation every remaining
spinor index is written dotted, the undotted $SU(2)$ index being carried by
the symbol $+$: the vector appears through the bispinor map
$(\sigma^m)_{+\dot\alpha}$ and $(\sigma^{mn})_{++}$ generates the undotted
(plus) $SU(2)$.
\begin{remark}\label{rem:2:pmframe}
Only the frame and the names $\nabla_\pm,\widetilde\nabla_{\dot\pm}$ are fixed
here. The connection-level definitions of the gauge-covariant derivatives are
deferred to Section~\ref{sec:gaugebv}, their use in the letter algebra to
Section~\ref{sec:superspace}.
\end{remark}

\subsection{Flat $N=1$ superspace}\label{sec:notation:superspace}

Flat $N=1$ superspace has coordinates
$z=(x^m,\theta^\alpha,\bar\theta_{\dot\alpha})$ (or $(x^\mu,\theta,\bar\theta)$
in Lorentzian signature); the odd coordinates are distinct from the
supersymmetry transformation parameters $\varepsilon,\bar\varepsilon$. We
write
\begin{equation}
\theta\sigma^m\bar\theta
:=\theta^\alpha(\sigma^m)_{\alpha\dot\beta}\bar\theta^{\dot\beta},\qquad
\theta\sigma^\mu\bar\theta
:=\theta^\alpha(\sigma^\mu)_{\alpha\dot\beta}\bar\theta^{\dot\beta} .
\label{eq:2:thetasigmabar}
\end{equation}
Odd derivatives are \emph{left} derivatives,
$\partial_\alpha:=\vec\partial/\partial\theta^\alpha$,
$\bar\partial^{\dot\alpha}:=\vec\partial/\partial\bar\theta_{\dot\alpha}$,
$\bar\partial_{\dot\alpha}
:=\epsilon_{\dot\alpha\dot\beta}\bar\partial^{\dot\beta}$, obeying the graded
Leibniz rule and
\begin{equation}
\{\partial_\alpha,\theta^\beta\}=\delta_\alpha{}^\beta,\qquad
\{\bar\partial^{\dot\alpha},\bar\theta_{\dot\beta}\}
=\delta^{\dot\alpha}{}_{\dot\beta},\qquad
\{\bar\partial_{\dot\alpha},\bar\theta^{\dot\beta}\}
=-\,\delta_{\dot\alpha}{}^{\dot\beta}.
\label{eq:2:grassmann}
\end{equation}
\begin{remark}\label{rem:2:barsign}
The minus sign in the last anticommutator of \eqref{eq:2:grassmann} is a
standard source of sign errors; it traces back to
$\bar\partial_{\dot\alpha}=\epsilon_{\dot\alpha\dot\beta}\bar\partial^{\dot\beta}$
with $\epsilon_{\dot1\dot2}=-1$.
\end{remark}
A supersymmetry transformation acts on the coordinates as
\begin{equation}
\theta'^\alpha=\theta^\alpha+\varepsilon^\alpha,\qquad
\bar\theta'_{\dot\alpha}=\bar\theta_{\dot\alpha}+\bar\varepsilon_{\dot\alpha},
\qquad
x'^\mu=x^\mu
+i\varepsilon^\alpha(\sigma^\mu)_{\alpha\dot\beta}\bar\theta^{\dot\beta}
-i\theta^\alpha(\sigma^\mu)_{\alpha\dot\beta}\bar\varepsilon^{\dot\beta} .
\label{eq:2:susytransf}
\end{equation}
We use the same symbols for the Noether charges and for their realizations as
differential operators on superspace, the difference being clear from context.
With $P_\mu=-i\partial_\mu$ the Lorentzian supercharges and supersymmetric
covariant derivatives are
\begin{equation}
Q_\alpha=-i\partial_\alpha
+(\sigma^\mu)_{\alpha\dot\beta}\bar\theta^{\dot\beta}\partial_\mu,\qquad
\bar Q_{\dot\alpha}=-i\bar\partial_{\dot\alpha}
-\theta^\beta(\sigma^\mu)_{\beta\dot\alpha}\partial_\mu ,
\label{eq:2:QL}
\end{equation}
\begin{equation}
\cD_\alpha=\partial_\alpha
-i(\sigma^\mu)_{\alpha\dot\beta}\bar\theta^{\dot\beta}\partial_\mu,\qquad
\bcD_{\dot\alpha}=\bar\partial_{\dot\alpha}
+i\theta^\beta(\sigma^\mu)_{\beta\dot\alpha}\partial_\mu ,
\label{eq:2:DL}
\end{equation}
where $\cD_\alpha,\bcD_{\dot\alpha}$ anticommute with all supercharges. Their
algebras are
\begin{equation}
\{Q_\alpha,\bar Q_{\dot\beta}\}
=2i(\sigma^\mu)_{\alpha\dot\beta}\partial_\mu
=-2(\sigma^\mu)_{\alpha\dot\beta}P_\mu,\qquad
\{\cD_\alpha,\bcD_{\dot\beta}\}
=2i(\sigma^\mu)_{\alpha\dot\beta}\partial_\mu
=-2(\sigma^\mu)_{\alpha\dot\beta}P_\mu ,
\label{eq:2:algebraL}
\end{equation}
with vanishing same-chirality anticommutators. The same relations hold for the
Hilbert-space Noether charges; the normalization is fixed by Hamiltonian
positivity, $\{Q_1,(Q_1)^\dagger\}+\{Q_2,(Q_2)^\dagger\}=4H$, and the central
charge vanishes. The Lorentzian chiral coordinates are
\begin{equation}
y^\mu:=x^\mu-i\theta\sigma^\mu\bar\theta,\qquad
\widetilde y^\mu:=x^\mu+i\theta\sigma^\mu\bar\theta \ (=\bar y^\mu),\qquad
\bcD_{\dot\alpha}y^\mu=0,\qquad
\cD_\alpha\widetilde y^\mu=0 .
\label{eq:2:ychiralL}
\end{equation}
After Wick rotation the coordinate momentum is $P_m=-\partial_m$ and
\begin{equation}
\begin{aligned}
Q_\alpha&=-\partial_\alpha
+(\sigma^m)_{\alpha\dot\beta}\bar\theta^{\dot\beta}\partial_m, &
\bar Q_{\dot\alpha}&=-\bar\partial_{\dot\alpha}
-\theta^\beta(\sigma^m)_{\beta\dot\alpha}\partial_m,\\
\cD_\alpha&=\partial_\alpha
+(\sigma^m)_{\alpha\dot\beta}\bar\theta^{\dot\beta}\partial_m, &
\bcD_{\dot\alpha}&=\bar\partial_{\dot\alpha}
-\theta^\beta(\sigma^m)_{\beta\dot\alpha}\partial_m ,
\end{aligned}
\label{eq:2:QDE}
\end{equation}
with algebras
\begin{equation}
\{Q_\alpha,\bar Q_{\dot\beta}\}
=+2(\sigma^m)_{\alpha\dot\beta}\partial_m
=-2(\sigma^m)_{\alpha\dot\beta}P_m,\qquad
\{\cD_\alpha,\bcD_{\dot\beta}\}
=-2(\sigma^m)_{\alpha\dot\beta}\partial_m
=+2(\sigma^m)_{\alpha\dot\beta}P_m .
\label{eq:2:algebraE}
\end{equation}
Same-chirality anticommutators vanish in both signatures, all mixed
$Q$-$\cD$ anticommutators vanish, and everything commutes with the momentum.
The Euclidean chiral coordinates are
\begin{equation}
y^m:=x^m+\theta\sigma^m\bar\theta,\qquad
\widetilde y^m:=x^m-\theta\sigma^m\bar\theta,\qquad
\bcD_{\dot\alpha}y^m=0,\qquad
\cD_\alpha\widetilde y^m=0 .
\label{eq:2:ychiralE}
\end{equation}
\begin{remark}[Sign trap]\label{rem:2:signtrap}
In Lorentzian signature the $Q$- and $\cD$-algebras coincide,
$\{Q,\bar Q\}=\{\cD,\bcD\}=-2\sigma^\mu P_\mu$; in Euclidean signature they
differ by a relative minus sign, $\{Q,\bar Q\}=-2\sigma^m P_m$ versus
$\{\cD,\bcD\}=+2\sigma^m P_m$. All sign conventions of the loop computation in
Sections~\ref{sec:superspace}--\ref{sec:component} are tied to
\eqref{eq:2:algebraE}.
\end{remark}

\subsection{Theta expansions and superfields}\label{sec:notation:superfields}

The basic monomials and their normalizations are
\begin{equation}
\theta^2:=\theta^\alpha\theta_\alpha=-2\theta^1\theta^2,\qquad
\bar\theta^2:=\bar\theta_{\dot\alpha}\bar\theta^{\dot\alpha}
=2\bar\theta_{\dot1}\bar\theta_{\dot2},\qquad
\theta^\alpha\theta^\beta=-\tfrac12\epsilon^{\alpha\beta}\theta^2,\qquad
\bar\theta_{\dot\alpha}\bar\theta_{\dot\beta}
=-\tfrac12\epsilon_{\dot\alpha\dot\beta}\bar\theta^2 .
\label{eq:2:theta2}
\end{equation}
With $\cD^2:=\cD^\alpha\cD_\alpha$ and
$\bcD^2:=\bcD_{\dot\alpha}\bcD^{\dot\alpha}$ one has, in either signature,
\begin{equation}
\cD^2\theta^2=-4,\qquad \bcD^2\bar\theta^2=-4 .
\label{eq:2:D2theta2}
\end{equation}
Writing $X|:=X|_{\theta=\bar\theta=0}$, the component projectors are
\begin{equation}
[X]_F:=-\tfrac14\cD^2X\big|,\qquad
[\widetilde X]_{\widetilde F}:=-\tfrac14\bcD^2\widetilde X\big|,\qquad
[Y]_D:=\tfrac1{16}\cD^2\bcD^2Y\big| \quad(\bcD^2\ \text{acts first}),
\label{eq:2:projectors}
\end{equation}
and the full superspace measure is
$\int d^4x\,d^4\theta\,Y:=\int d^4x\,[Y]_D$.
\begin{remark}\label{rem:2:generalsf}
A general (unconstrained) superfield enters this paper only through the
projectors \eqref{eq:2:projectors}; its full fourteen-component expansion is
never used. The vector multiplet is always reduced to Wess--Zumino gauge,
\eqref{eq:2:wzE} below.
\end{remark}
A chiral superfield $\Phi$ and an antichiral superfield $\widetilde\Phi$ obey
\begin{equation}
\bcD_{\dot\alpha}\Phi=0,\qquad \cD_\alpha\widetilde\Phi=0 ,
\label{eq:2:chirality}
\end{equation}
and have components $(\phi,\psi_\alpha,F)$ and
$(\widetilde\phi,\widetilde\psi_{\dot\alpha},\widetilde F)$ given by
\begin{equation}
\phi:=\Phi\big|,\qquad
\psi_\alpha:=\tfrac1{\sqrt2}\cD_\alpha\Phi\big|,\qquad
F:=-\tfrac14\cD^2\Phi\big|;\qquad
\widetilde\phi:=\widetilde\Phi\big|,\qquad
\widetilde\psi_{\dot\alpha}
:=\tfrac1{\sqrt2}\bcD_{\dot\alpha}\widetilde\Phi\big|,\qquad
\widetilde F:=-\tfrac14\bcD^2\widetilde\Phi\big| .
\label{eq:2:chiralcomponents}
\end{equation}
In Lorentzian signature $\widetilde\Phi=\bar\Phi$ and
$(\widetilde\phi,\widetilde\psi,\widetilde F)=(\bar\phi,\bar\psi,\bar F)$; in
Euclidean signature $\widetilde\Phi$ is independent of $\Phi$
(Remark~\ref{rem:2:noconjugation}). Their chiral-coordinate expansions are
\begin{equation}
\Phi(y,\theta)=\phi(y)+\sqrt2\,\theta^\alpha\psi_\alpha(y)+\theta^2F(y),
\qquad
\widetilde\Phi(\widetilde y,\bar\theta)
=\widetilde\phi(\widetilde y)
+\sqrt2\,\bar\theta_{\dot\alpha}\widetilde\psi^{\dot\alpha}(\widetilde y)
+\bar\theta^2\widetilde F(\widetilde y),
\label{eq:2:chiralexpansion}
\end{equation}
with $y,\widetilde y$ from \eqref{eq:2:ychiralL} or \eqref{eq:2:ychiralE}.
Expanding in $x$-space gives, in Euclidean signature,
\begin{equation}
\Phi=\phi+\sqrt2\,\theta\psi+\theta^2F
+\theta\sigma^m\bar\theta\,\partial_m\phi
-\frac1{\sqrt2}\,\theta^2
(\partial_m\psi^\alpha)(\sigma^m)_{\alpha\dot\beta}\bar\theta^{\dot\beta}
+\frac14\,\theta^2\bar\theta^2\,\partial^m\partial_m\phi ,
\label{eq:2:chiralxE}
\end{equation}
\begin{equation}
\widetilde\Phi=\widetilde\phi
+\sqrt2\,\bar\theta\widetilde\psi
+\bar\theta^2\widetilde F
-\theta\sigma^m\bar\theta\,\partial_m\widetilde\phi
+\frac1{\sqrt2}\,\theta^\alpha(\sigma^m)_{\alpha\dot\beta}
(\partial_m\widetilde\psi^{\dot\beta})\,\bar\theta^2
+\frac14\,\theta^2\bar\theta^2\,\partial^m\partial_m\widetilde\phi ,
\label{eq:2:antichiralxE}
\end{equation}
where $\partial^m\partial_m$ is the Euclidean Laplacian; the Lorentzian
counterparts follow from \eqref{eq:2:chiralexpansion} with
\eqref{eq:2:ychiralL}.

The vector multiplet is described by the prepotential $\cV=\cV^AT_A$ with
$[T_A,T_B]=if_{AB}{}^CT_C$ in an orthonormal basis,
$\kappa_{AB}=\delta_{AB}$. Wess--Zumino gauge is the residual gauge choice
\begin{equation}
\cV\big|=\theta\cV\big|=\bar\theta\cV\big|=0 ,
\label{eq:2:wzcondition}
\end{equation}
in which the prepotential reduces to
\begin{equation}
\cV_{\rm WZ}=-2\theta\sigma^\mu\bar\theta\,A_\mu
+2i\theta^2\bar\theta_{\dot\alpha}\bar\lambda^{\dot\alpha}
-2i\bar\theta^2\theta^\alpha\lambda_\alpha
+\theta^2\bar\theta^2 D
\qquad\text{(Lorentzian)},
\label{eq:2:wzL}
\end{equation}
\begin{equation}
\cV_{\rm WZ}=-2i\theta\sigma^m\bar\theta\,A_m
+2i\theta^2\bar\theta_{\dot\alpha}\widetilde\lambda^{\dot\alpha}
-2i\bar\theta^2\theta^\alpha\lambda_\alpha
+\theta^2\bar\theta^2 D
\qquad\text{(Euclidean)},
\label{eq:2:wzE}
\end{equation}
with component fields $(A_m,\lambda_\alpha,\widetilde\lambda_{\dot\alpha},D)$;
note that the Euclidean prepotential contains the \emph{independent}
antichiral gaugino $\widetilde\lambda_{\dot\alpha}$. The gauge field is the
ordered-commutator projection
\begin{equation}
A_\mu=\frac18(\bar\sigma_\mu)^{\dot\beta\alpha}
[\cD_\alpha,\bcD_{\dot\beta}]_{\rm ord}\cV\Big|
\quad\text{(Lorentzian)},\qquad
A_m=\frac{i}{8}(\bar\sigma_m)^{\dot\beta\alpha}
[\cD_\alpha,\bcD_{\dot\beta}]_{\rm ord}\cV\Big|
\quad\text{(Euclidean)},
\label{eq:2:Aprojection}
\end{equation}
where $[\cdot,\cdot]_{\rm ord}$ keeps the spinor derivatives in the displayed
order. The gaugino field-strength superfields are
\begin{equation}
\cW_\alpha:=-\frac18\,\bcD^2\big[e^{-\cV}(\cD_\alpha e^{\cV})\big],\qquad
\bcW_{\dot\alpha}
:=+\frac18\,\cD^2\big[e^{\cV}(\bcD_{\dot\alpha}e^{-\cV})\big],
\label{eq:2:fieldstrengths}
\end{equation}
chiral and antichiral respectively, $\bcD_{\dot\alpha}\cW_\beta=0$,
$\cD_\alpha\bcW_{\dot\beta}=0$, with component projections
\begin{equation}
\lambda_\alpha:=i\cW_\alpha\big|,\qquad
\widetilde\lambda_{\dot\alpha}:=-i\bcW_{\dot\alpha}\big|,\qquad
D:=-\frac12\cD^\alpha\cW_\alpha\big|
=+\frac12\bcD^{\dot\alpha}\bcW_{\dot\alpha}\big| .
\label{eq:2:vectorcomponents}
\end{equation}
In Lorentzian signature $\bcW_{\dot\alpha}=\bar\cW_{\dot\alpha}$ and
$\widetilde\lambda_{\dot\alpha}=\bar\lambda_{\dot\alpha}$; in Euclidean
signature $(\cW_\alpha,\bcW_{\dot\alpha})$ and
$(\lambda_\alpha,\widetilde\lambda_{\dot\alpha})$ are independent.

\subsection{Chiral and vector representations}\label{sec:notation:representations}

Gauge transformations are generated by a chiral parameter $\Lambda$ and an
independent antichiral parameter $\bar\Lambda$,
\begin{equation}
h:=e^{i\Lambda},\qquad \bar h:=e^{i\bar\Lambda},\qquad
\bcD_{\dot\alpha}\Lambda=0,\qquad \cD_\alpha\bar\Lambda=0 .
\label{eq:2:gaugeparams}
\end{equation}
In the \emph{chiral representation} superfields are functions of the chiral
coordinates \eqref{eq:2:ychiralL}--\eqref{eq:2:ychiralE}, and
\begin{equation}
\Phi'=h\,\Phi,\qquad
\widetilde\Phi'=\widetilde\Phi\,\bar h^{-1},\qquad
(e^{\cV})'=\bar h\,e^{\cV}\,h^{-1} .
\label{eq:2:chiralframe}
\end{equation}
Thus the bridge $e^{\cV}$ carries the entire gauge field in this frame: the
combination $\widetilde\Phi\,e^{\cV}\,\Phi$ is gauge invariant, the gauge
connection appears only through $e^{-\cV}(\cD_\alpha e^{\cV})$ in the undotted
spinor derivative, and the field strengths \eqref{eq:2:fieldstrengths}
transform homogeneously,
\begin{equation}
\cW'_\alpha=h\,\cW_\alpha\,h^{-1},\qquad
\bcW'_{\dot\alpha}=\bar h\,\bcW_{\dot\alpha}\,\bar h^{-1} .
\label{eq:2:Wtransform}
\end{equation}
In the \emph{vector representation} one works instead on the full
$(x,\theta,\bar\theta)$ superspace: all superspace derivatives carry
connections, the matter fields are covariantly chiral, and gauge
transformations act by a single real group element, so that the complexified
parameters $h,\bar h$ and the explicit bridge disappear from the
transformation laws. The two representations are related by a factorization of
the bridge $e^{\cV}$ and are physically equivalent; we use the chiral
representation in what follows.
\begin{remark}\label{rem:2:frames}
This subsection states the two frames at the superfield level only. The
connection-level definitions of the gauge-covariant derivatives
$\nabla_\alpha,\widetilde\nabla_{\dot\alpha}$ of \eqref{eq:2:pmframe}, the
bridge factorization, the gauge action, and the BV--BRST formulation are
deferred to Section~\ref{sec:gaugebv}.
\end{remark}
